\section{我们在 AI Bubble 中吗：一个仪表盘而不是一句口号}

上一章把收入池拆清楚，本章回到节目标题里的泡沫问题。Freda 的回答是：今天不是 bubble，但未来可能会变成 bubble。她反对空谈情绪，主张看两个层面：第一，今天是否已经泡沫；第二，未来是否会吹成泡沫。判断依据包括 AI 应用落地速度、收入兑现、大公司 ROIC、模型是否继续进步，以及 AI 收入是否持续增加。

\begin{figure}[H]
\centering
\includegraphics[width=0.96\textwidth]{ai-bubble-dashboard.png}
\caption{判断 AI Bubble 的仪表盘：用户、收入、ROIC、模型进步和估值。自制概念图，依据 01:14:11--01:16:31 对谈内容整理。}
\end{figure}

\begin{knowledgebox}{读图：先看可验证指标}
用户渗透速度说明产品是否真实普及；OpenAI/Anthropic 收入说明商业化是否兑现；大厂 ROIC 说明投入是否有回报；模型是否继续 scale 决定长期故事；估值则要和这些指标一起看，不能单独作为泡沫证据。
\end{knowledgebox}

\subsection{为什么她认为今天还不是 bubble}

本节先看“今天”这个时间点。前面已经把判断指标列成仪表盘，现在要问的是：这些指标是否已经显示纯叙事透支？

节目给出的理由包括：GPT 出现不到三年，用户普及程度已经相当于互联网发展十年的成果；AI 已经有几百亿美元收入；大公司投入 AI 后季度 ROIC 仍在提升；OpenAI 和 Anthropic 是可跟踪的大头，Coding、视频生成、音频、客服、法律、医疗等垂类也已经出现收入。也就是说，今天的 AI 至少不是纯叙事，已经有用户、收入和财务效果。

\subsection{未来会不会变成 bubble}

接下来把问题从现在推到未来。即使今天不是泡沫，也可能在收入、模型进步或估值三者脱节时变成泡沫。

未来是否泡沫，Freda 看两个变量。第一，模型是否继续进步。Gemini 3 的出现被她视为 pre-training 仍可继续 scale 的信号，Blackwell 训练出的下一代模型也值得观察。第二，AI 收入是否继续增加。如果收入跟不上 CAPEX、估值和市场预期，泡沫风险就会上升。

\begin{warningbox}{不是 bubble 不等于没有泡沫公司}
一个大周期真实存在，不代表所有公司估值合理，也不代表每个赛道都会成功。铁路、互联网和云计算都是真革命，但每轮也都有大量失败公司和过度资本开支。
\end{warningbox}

\subsection{本章小结}

泡沫判断最好变成仪表盘：看用户、收入、ROIC、模型进步和估值，而不是只看市场情绪。当前 AI 已经有真实收入和应用，但未来是否泡沫取决于收入能否持续兑现、模型能否继续进步。

